\documentclass[10pt,a4paper]{amsart}
\usepackage[margin=19mm]{geometry}
\usepackage{amsmath,amssymb,mathtools}
\usepackage[scr=rsfs,cal=euler]{mathalfa}
\usepackage{xfrac}
\usepackage[unicode,hidelinks,bookmarks=false]{hyperref}
\newcommand{\ie}{i.e.\ }
\newcommand{\cf}{cf.\ }
\newcommand{\resp}{respectively\ }
\newcommand{\Subsection}{\subsection}
\providecommand{\nobreakdash}{\nobreak\hbox{-}}
\setlength{\emergencystretch}{2em}
\multlinegap0pt
\usepackage{bm}
\usepackage{bbm}
\usepackage{dsfont}
\usepackage{xspace}
\usepackage{cancel}
\usepackage{mathtools}
\usepackage{amsthm}
\makeatletter
\@ifundefined{theorem}{\newtheorem{theorem}{Theorem}}{}
\@ifundefined{lemma}{\newtheorem{lemma}[theorem]{Lemma}}{}
\makeatother
\def\P{{\mathds{P}}}
\def\R{{\mathds{R}}}
\def\bd{{\partial}}
\def\cC{{\mathcal C}}
\def\cF{{\mathcal F}}
\def\cH{{\mathcal H}}
\def\cX{{\mathcal X}}
\def\r{{\rho}}
\newcommand{\IN}{\mathds{N}}
\newcommand{\diam}{\mathrm{diam}}
\title[Central limit theorems for random boundary polytopes]{Central limit theorems for random boundary polytopes}
\author{Matthias Reitzner, Mathias Sonnleitner}
\date{Source: arXiv:2509.20058v1 (CC BY 4.0)}
\begin{document}
\maketitle

\noindent\emph{Source and licence.} Central limit theorems for random boundary polytopes. Version arXiv:2509.20058v1 (CC BY 4.0), distributed under the Creative Commons Attribution 4.0 International licence, as recorded in the arXiv licence metadata for this exact version and shown on the abstract page https://arxiv.org/abs/2509.20058v1. Full rights notice: licenses/arxiv-2509.20058-CC-BY-4.0.txt.\par\smallskip

\noindent\textit{Editorial scope.} Counted unit: the Theorem whose source label is \texttt{th:rad\_influence} in the article (source0.tex lines 551--558, source label \texttt{th:rad\_influence}), delivered together with the article's own complete original proof of it (source0.tex lines 559--641, one \texttt{proof} environment of 83 lines). No mathematics is written, repaired or re-derived: statement and proof are the article's own text line for line. Required context retained uncounted: lem:cap-convex (lines 310--324); eq:in-out (lines 327--331). Every definition, display and label of those lines is retained unchanged. Omitted from this package and not redistributed: 1--309, 325--326, 332--550, 642--736. Nothing inside a retained excerpt is omitted or altered: every retained body is delivered exactly as the article's lines are written, so no omission receipt of any kind accompanies this package. Citations: the retained excerpts carry no citation command at all, so no bibliography entry of the article is transcribed and no reference is redistributed. Reference adaptation: none was needed and none was made; every reference inside the delivered unit resolves inside it, so no unresolved reference or citation is produced. Printed slips kept literal: no printed word, symbol, hypothesis or number of the counted statement or of its proof was altered. Layout and class: the article's own class is \texttt{article} with packages including geometry, graphicx, inputenc, color, latexsym, amsbsy, amsthm, amssymb, amsmath, amsfonts, dsfont, enumerate, hyperref; the wrapper is the lane's pinned \texttt{amsart} editorial wrapper at 10pt on A4, which restates the theorem-like environments the retained text needs and adds the article's own macro definitions that the retained text needs (\texttt{\textbackslash IN}, \texttt{\textbackslash P}, \texttt{\textbackslash R}, \texttt{\textbackslash bd}, \texttt{\textbackslash cC}, \texttt{\textbackslash cF}, \texttt{\textbackslash cH}, \texttt{\textbackslash cX}, \texttt{\textbackslash diam}, \texttt{\textbackslash r}), with extra wrapper packages (bm, bbm, dsfont, xspace, cancel, mathtools). The delivered numbering is the unit's own. Assets: no retained span contains a figure environment, a graphics-inclusion command, a PDF-page inclusion, a table or a TikZ picture; no image file is copied into this package, the package declares no asset dependency and no image file is redistributed with it. Licence: version arXiv:2509.20058v1 of this article is distributed under the Creative Commons Attribution 4.0 International licence, as recorded in the arXiv licence metadata for this exact version and shown on the abstract page https://arxiv.org/abs/2509.20058v1. A full rights notice is delivered at licenses/arxiv-2509.20058-CC-BY-4.0.txt. Characters: every retained excerpt is pure ASCII, so the delivered engine, XeLaTeX with its default Latin Modern text font, reports no missing character.
\begin{lemma}\label{lem:cap-convex}
Let $K\in \cC_+^2$. Then, for any $x\in \bd K$ and $h\le \diam(K)$, 
\[
c\, h^{\frac{d-1}{2}}
\le \cH^{d-1}\big(C^{\bd K}(x,h)\big)
\le C\, h^{\frac{d-1}{2}},
\]
and for $r\le \diam(K)$,
\begin{equation*} 
c \, r^{d-1}
\le \cH^{d-1}\big(B(x,r)\cap \bd K\big)
\le C\, r^{d-1},
\end{equation*}
where $B(x,r)$ denotes a Euclidean ball with center $x$ and radius $r$, $\diam(K)$ is the diameter of $K$, and $c,C>0$ only depend on $K$.
\end{lemma}

\begin{equation} \label{eq:in-out}
B(x-r_{in} u_x, r_{in}) 
\subset K
\subset B(x-r_{out} u_x, r_{out}).
\end{equation}

\begin{theorem}\label{th:rad_influence}
There exist $c,C>0$ such that for each $x \in \bd K$ and $r>0$ we have
$$
\P( R(x, \cX_{n} \cup\{x\}) \geq r)
\leq
C \exp(- c r^{d-1} n),\qquad n\in \IN. 
$$
\end{theorem}

\begin{proof}
	Recall Blaschke's rolling theorem in the form \eqref{eq:in-out} with a ball of radius $r_{in}$ rolling inside and a ball of radius $r_{out}$ rolling outside. Since $R\le {\rm diam}(K)$, it is sufficient to show the statement of Theorem~\ref{th:rad_influence} for $r\le r_{in}$.

	Let $x\in \bd K$ and $r\le r_{in}$. If $R(x, \cX_n\cup \{x\}) \geq r$, then there is a hyperplane $H_F$ containing a facet $F \in \cF_{d-1}([\cX_n,x])$ with $x$ as one of its vertices such that the diameter of $H_F \cap K$ exceeds $r$. Thus, also the diameter of $H_F \cap B(x-r_{out} u_x, r_{out})$ exceeds $r$. Hence the diameter of the ball $H_F \cap B(x-r_{in} u_x, r_{in})$ exceeds $\frac{r_{in}}{r_{out}} r$, and its radius exceeds $\frac{r_{in}}{2r_{out}} r$. 
Because of
$$ x \in H_F \cap B(x-r_{in} u_x, r_{in}) \subset H_F \cap K , $$
there is a ball with center $\tilde x \in \bd K$, of radius $\frac{r_{in}}{2r_{out}} r$, which contains $x$ in its boundary, such that 
\begin{equation}\label{eq:cap_ball}
B\left(\tilde x, \frac{r_{in}}{2r_{out}} r \right) \cap \bd K \subset H_F^+ \cap \bd K 
.
\end{equation}
We use the following geometric lemma.
\begin{lemma}\label{le:points_cover}
For any $x  \in \bd K$ there are $z_1, \dots , z_{2^{d-1}} \in \bd K$ such that the following holds:
if $\tilde x \in \bd K$ with $\|\tilde x -  x\| = \rho < r_{in}$, then there is an $i \in \{1, \dots , 2^{d-1}\}$ such that 
$$ B\left(z_i, \frac 1{4d^2} \rho \right) \subset {\rm int}\, B(\tilde x, \rho) . $$ 
\end{lemma}
\begin{proof}
Without loss of generality assume that $x$ is the origin, and that the outer unit normal vector of $x$ equals $-e_d$. The orthogonal hyperplanes to the coordinate vectors $e_1, \dots, e_{d}$ dissect $\R^d$ into $2^d$ orthants from which those which contain $e_d$ may have non-empty intersection with the interior of $K$, say $O_1, \dots,  O_{2^{d-1}}$.
Assume that $O_1 = \{x \in \R^d \colon x_1, \dots, x_d \geq 0 \}$, and that $\tilde x \in O_1$. 

We define
$$z_1 = \left(\frac \r {\sqrt {d(d-1)}}, \dots, \frac \rho {\sqrt {d(d-1)}} , h_d \right)    , $$
where $h_d\ge 0$ is chosen such that $z_1 \in \bd K $ (and $z_2, \dots , z_{2^{d-1}}$ analogously).
Because $\tilde x$ and $z_1$ are not in the interior of the ball with radius $r_{in}$, we have
$$
h_d \leq \frac {\r^2}{d r_{in}} \leq \frac {\r}{d }
, \text{ and }\ 
\tilde x_d \leq \frac {\rho^2}{2 r_{in}} \leq \frac {\rho}{2} 
$$
We estimate the distance between $z_1$ and $\tilde x$.
\begin{align*}
\| \tilde x -z_1  \|^2
&=
\| \tilde x \| ^2  -  2 \tilde x \cdot  z_{1}  +  \|z_{1}\|^2
\\ & \leq 
\rho^2   -   \frac{2\rho}{\sqrt{d(d-1)}} \sum_{i=1}^d \tilde x_i + \frac{2 \rho}{\sqrt{d(d-1)}} \tilde x_d - 2 h_d \tilde x_d  +  \rho^2 \frac 1d  + h_d^2
\\ & \leq
\rho^2 \left(1 + \frac 1 d - \frac{2}{\sqrt{d(d-1)}}\right) + \frac{2\rho }{\sqrt{d(d-1)}} \tilde x_d - 2 \tilde x_d h_d  + h_d^2
\end{align*}
because $\sum _1^d \tilde x_i \geq \| \tilde x \|_2 = \rho$ for $\tilde x_i \geq 0$.
The function $ \frac{2\rho }{\sqrt{d(d-1)}} \tilde x_d - 2 \tilde x_d h_d  + h_d^2 $ has no maximum in the interior of the region 
$(\tilde x_d, h_d) \in[0, \frac \rho 2] \times[0, \frac \rho d]$, and attains for $\tilde x_d= \frac \rho 2$, $h_d = 0$ its maximum $\frac{\rho ^2}{\sqrt{d(d-1)}}$. Hence
\begin{align*}
\| \tilde x -z_1  \|^2
& \leq
\rho^2 \left(1 + \frac 1 d - \frac{1}{\sqrt{d(d-1)}} \right)\\
& \leq  
\rho^2 \left(1 - \frac{1}{2d^2} \right)
\end{align*}
since $(1-\frac{1}{d})^{-1/2}\ge 1+\frac{1}{2d}$.
This proves that the ball with center $z_1$ of radius
$$
\rho- \sqrt{ \rho^2 - \frac{\rho^2}{2d^2}  } > \frac 1{4d^2} \rho
$$
is contained in the interior of $B(\tilde x, \rho)$. This completes the proof of Lemma~\ref{le:points_cover}.
\end{proof}

From \eqref{eq:cap_ball} and Lemma \ref{le:points_cover} we deduce that for each boundary point $x$ and for any fixed $r$ there are points $z_1, \dots, z_{2^{d-1}}$ such that any cap of diameter at least $r$ contains one of the balls
$$
B \left(z_i, \frac{r_{in}}{8 d^2 r_{out}} r \right),\qquad i=1,\dots,2^{d-1}.
$$
This implies
\begin{align*}
\P( R(x, \cX_{n} \cap\{x\}) \geq r)
& =
\P\big( \exists F \in \cF([x,\cX_{n}])\colon x \in F,\, {\rm diam}( H_F\cap K ) \geq r \big)
\\ & \leq
\sum_{i=1}^{2^{d-1}} \P  \left( B\left(z_i, \frac{r_{in}}{8 d^2 r_{out}} r \right) \cap \cX_{n} = \emptyset \right)
\\ & \leq
\sum_{i=1}^{2^{d-1}} \left(1- \cH^{d-1} \left(B\left(z_i, \frac{r_{in}}{8 d^2 r_{out}} r \right) \cap \bd K\right) \right)^{n}.
\end{align*}

With $c'=(\frac{r_{in}}{8 d^2 r_{out}})^{d-1}c$, where $c$ is from Lemma \ref{lem:cap-convex}, it holds that
\begin{align*}
\P( R(x, \cX_{n} \cap\{x\}) \geq r)
& \leq
2^{d-1} \left(1- c' r^{d-1} \right)^{n} 
\\ & \leq
2^{d-1} \exp(- c' r^{d-1} n).
\end{align*}
This completes the proof of Theorem~\ref{th:rad_influence}.
\end{proof}

\end{document}
